\section{Invariant Euler nonvanishing by a fixed invariant system}\label{har:sec:eulerroute}

We now give a harmonic proof of Corollary~\ref{thm:main}. Starting with
the determinant sections, we choose the base from powers of one fixed
invariant line. Maximality of its image will force the invariant
adjoint spaces on the generic fiber to have rank one. The determinant
intersection identity will supply the curvature-kernel condition.
Together these properties give the continuously semipositive base line
and the vanishing at twist zero required by
Proposition~\ref{har:prop:euler}. The premise here is the nonzero
invariant Euler characteristic; vanishing of every higher
structure-sheaf cohomology group is not required.

\begin{proof}[Harmonic proof of Corollary~\ref{thm:main}]
Use Lemma~\ref{har:lem:forms}, retaining the cases which did not already
produce a section, and write \(N_i=M+a_iL\), \(s_i\) for its invariant
sections. Write \(n=\dim X\), let \(\theta\) be the averaged smooth
semipositive curvature, let \(r\) be its maximal rank, and fix a
K\"ahler form \(\omega\). Choose indices \(a_-<a\) and keep unbounded indices above
\(a\). Put \(N=M+aL\). Among rational maps given by all invariant
sections of positive powers of \(N\), choose one with maximal image
dimension. Resolve its base ideal equivariantly and take Stein
factorization. We obtain a smooth projective \(X_0\), a birational
\(\pi_0:X_0\to X\), and \(f_0:X_0\to Y_0\) with connected fibers and
normal projective base, with
\begin{equation}\label{har:eq:iitaka}
 d\pi_0^*N=f_0^*H_0+G_0,\qquad G_0\ge0,
\end{equation}
where \(H_0\) is ample. The torus acts trivially on \(Y_0,H_0\), and
the defining section of \(G_0\) is invariant.

The invariant sections of \(j\pi_0^*N\) on the generic fiber form a
one-dimensional vector space over \(\mathbb C(Y_0)\) for every \(j>0\).
Indeed finitely many generic-fiber sections extend rationally with
vertical poles. A section of a sufficiently high power of \(H_0\)
clears those poles; multiplication by the corresponding power of the
section of \(G_0\) gives invariant sections of a higher power of
\(\pi_0^*N\), and hence of \(N\) on \(X\). If two original sections
were independent, their ratio would not lie in \(\mathbb C(Y_0)\).
This field is relatively algebraically closed by Stein factorization,
so that ratio would be transcendental over it. Combine it with the
original image coordinates, by products to put all ratios in one
power. The image dimension increases, contradicting maximality.

For each \(a_i>a\) there are positive integers \(p,q\) such that
\(pa_i+qa_-=(p+q)a\). The invariant section
\(s_i^p s_-^q\) is thus generically proportional to \(s_a^{p+q}\).
Its effective generic-fiber divisor shows that all horizontal
components of \(\operatorname{div}(\pi_0^*s_i)\) lie in the fixed finite
support of \(\operatorname{div}(\pi_0^*s_a)\). Extract a subsequence of
their integer multiplicity vectors that is componentwise nondecreasing.
For two indices \(a_j>a_i\) in it, put
\[
 b=a_j-a_i>0,\qquad s=\pi_0^*s_j/\pi_0^*s_i.
\]
This is an invariant rational section of \(b\pi_0^*L\), regular on the
generic fiber. The relative canonical divisor
\(E_0=K_{X_0/X}\) has its horizontal support in the same fixed support:
\(\pi_0\) can be chosen to be an isomorphism outside the invariant base
ideal, which lies in the zero locus of \(s_a\). Fix \(k\ge0\), and work on the generic fiber.  The effective
divisors of \(s_{E_0}\) and \(s\) have support in the divisor of
\(\pi_0^*s_a\).  All coefficients are finite, so a sufficiently large
integer \(h\) satisfies
\[
 h\operatorname{div}(\pi_0^*s_a)
 \ge \operatorname{div}(s_{E_0})+k\operatorname{div}(s).
\]
The quotient \((\pi_0^*s_a)^h/(s_{E_0}s^k)\) is therefore a
regular invariant section of the difference line.  Multiplication
by it injects the entire invariant adjoint section space into the
one-dimensional invariant space for \(h\pi_0^*N\).  Consequently
\begin{equation}\label{har:eq:invrank}
 \dim_{\mathbb C(Y_0)}H^0((X_0)_\eta,E_0+kb\pi_0^*L)^T=1
 \qquad(k\ge0).
\end{equation}
If \(\dim Y_0=0\), the section \(s\) is already regular on \(X_0\)
and descends to the required section on \(X\).

Otherwise apply Lemma~\ref{val:flat-model} to this fibration and resolve
its normal equidimensional model:
\[
 V\longrightarrow W\longrightarrow X_0,\qquad g:W\to Y,
 \qquad f:V\to Y.
\]
Write \(\mu:W\to X\) and \(\pi:V\to X\) for the composed birational maps. We also use \(f_0\) for its pullback to \(V\). Here \(Y\) is smooth projective, \(W\) is normal equidimensional over
\(Y\), and \(V\) is smooth. All constructions are equivariant, with
trivial base action. The ranks
\eqref{har:eq:invrank} persist: a further effective birational canonical
divisor has trivial pushforward also on the generic smooth fiber.

It remains to relate this algebraic fibration to the curvature. The
kernel condition for \(D_V=\pi^*L\) follows from
Lemma~\ref{har:lem:null-divisors} and \eqref{har:eq:iitaka}. If \(r<n\),
\[
 \int_V (\pi^*\theta)^r\wedge f_0^*\eta_{H_0}
                       \wedge(\pi^*\omega)^{n-r-1}=0.
\]
Indeed the whole class \(dN\) has zero intersection, and both the
moving and effective fixed parts have nonnegative intersection. On
an ordinary open where \(\pi\) is an isomorphism, \(\theta\) has rank
\(r\), and the base-to-image map has injective differential, positivity
forces the horizontal form to annihilate \(\ker\pi^*\theta\). Its
kernel there is \(\ker df\). Thus
\(\ker\pi^*\theta\subset\ker df\). These algebraic generic conditions
meet the nonempty ordinary maximal-rank open. For full curvature rank
the kernel inclusion holds directly.

Proposition~\ref{har:prop:maximum} normalizes \(s\) to
\(g^*Q\to b\mu^*L\), with effective zero divisor \(G\), and gives a
continuous semipositive metric on \(Q\). Theorem~\ref{har:thm:kernel}
and Proposition~\ref{har:prop:euler} now produce an invariant section
of a positive power of \(L\).

\end{proof}

\subsection{An alternative proof that the maximum stays positive}

The preceding proof is complete. We record a second argument for the
positivity of the fiberwise maximum on its model. Instead of extending
the maximum weight across a codimension-two set, we first prove
nefness by stabilizing the sheaves of allowed poles. Blowing up a
hypothetical zero fiber then contradicts nefness. This argument uses
harmonic vanishing and is independent of the maximum-principle step
in Proposition~\ref{har:prop:maximum}.

\begin{proposition}[Pole-sheaf stabilization on the invariant Iitaka model]\label{har:prop:pole-stabilization}
Let $V\to W\to X$, $g:W\to Y$, $f:V\to Y$, and $D_V=\pi^*(-K_X)$ be the normal equidimensional model and smooth resolution in the preceding invariant Iitaka construction. Put $E=K_{V/X}$. Retain its invariant adjoint rank-one condition \eqref{har:eq:invrank} and its curvature-kernel inclusion, and write $bD_V=G+f^*Q$ after vertical normalization, with $G$ pulled back from the normal model and omitting a divisorial lift of each base prime. Then the increasing invariant pole sheaves
$\mathcal P_k=(f_*\mathcal O_V(E+kG))^T$, viewed as rational functions through $s_Es_G^k$, stabilize. The line $Q$ is nef and $G$ contains no entire fiber. These conclusions give a second proof that the fiberwise maximum defines a positive continuous metric.
\end{proposition}
\begin{proof}
On \(V\), form
\[
 \mathcal P_k=(f_*\mathcal O_V(E+kG))^T\qquad(k\ge0),
\]
viewed as sheaves of rational functions on \(Y\) by division by
\(s_Es_G^k\). They have rank one, contain \(\mathcal O_Y\), and increase
with \(k\). For each prime \(P\) of \(Y\), take a lift \(\Gamma\)
omitted by \(G\). Regularity of \(f^*a\,s_Es_G^k\) gives
\[
 \operatorname{ord}_{\Gamma}(f^*P)\operatorname{ord}_P(a)
      +\operatorname{ord}_{\Gamma}E\ge0.
\]
The possible poles are therefore bounded independently of \(k\) by
one divisor \(C\). Normality embeds all \(\mathcal P_k\) into
\(\mathcal O_Y(C)\), so Noetherianity gives a stable value
\(\mathcal P_\infty\).

Take a very ample \(H\) on \(Y\). Apply
Theorem~\ref{har:thm:kernel}, with trivial base twist, to the smooth coefficient
\((1+kb)D_V+lf^*H\), \(k,l\ge0\). Its curvature dominates the original
curvature, so its kernel is still vertical on the same open set.
Since its adjoint line is \(E+kG+f^*(kQ+lH)\), projection to
invariants and the projection formula give
\[
 H^i(Y,\mathcal P_k\otimes(kQ+lH))=0\qquad(i>0).
\]
For large \(k\), Castelnuovo--Mumford regularity makes
\(\mathcal P_\infty\otimes(kQ+(\dim Y)H)\) globally generated. Pull
back to the normalization of any curve and remove torsion. The fixed
sheaf has positive rank, so its degree plus \(k\) times that rank times
\(\deg Q\) is nonnegative for unbounded \(k\). Thus \(Q\) is nef.

If \(G\) contained the entire fiber over \(y\in Y\), blow up \(y\)
(the curve case is already excluded by normalization). The normalized
main transform remains equidimensional because it lies in a base
change with the same fiber-dimension upper bound. Recompute the
vertical minima on this model. The new rational base divisor is the
pullback of \(Q\) plus a strictly positive multiple of the exceptional
divisor: every divisorial lift of its generic point lies over the
old fiber and therefore has positive order in the pulled-back \(G\).
The rank and kernel arguments just given still apply, so the new
divisor is nef after clearing denominators. Its degree on a line in
the exceptional projective space is negative, a contradiction. Hence
no entire fiber lies in \(G\). Properness and openness give a positive
continuous maximum, and the moment calculation
\eqref{har:eq:moments} proves semipositivity without the codimension-two
extension argument. This gives the same conclusion by a distinct
pole-sheaf stabilization proof.
\end{proof}
